% WHAT AN \endinput IN A FILE NAME ENDS.
%
% \endinput is EXPANDABLE, so one written where a file name is still being
% scanned takes hold THERE rather than waiting to be obeyed -- and what it
% takes hold of is whatever file is being read when the next LINE is
% wanted. For
%
%   \input zzendin\endinput\relax
%
% that is zzendin, which comes out ONE LINE LONG: its first line is read
% and no more. One space further on --
%
%   \input zzendin \endinput
%
% -- the name has already ended, so the \endinput is read after zzendin has
% been read out, and ends the file that wrote it instead.
%
% At \tracingcommands=2 the \endinput draws a line of its own where the name
% is being scanned, because that is where it is expanded.
%
% HSTeX dropped an \endinput that ended a file name, on the ground that
% what is backed up after a name is not where one can take hold. The first
% half of that is right and the second is not: it takes hold of the file
% that has not been opened yet. trip line 424 writes
% `\expandafter\stopinput\input tripos\endinput\input', and read three
% lines of tripos where the reference reads one.
%
% This probe writes a file of its own; run it where that is welcome.
\catcode`\{=1 \catcode`\}=2
\tracingonline=1 \tracingcommands=2
\immediate\openout9=zzendin
\immediate\write9{\noexpand\message{[line one]}}
\immediate\write9{\noexpand\message{[line two]}}
\immediate\write9{\noexpand\message{[line three]}}
\immediate\closeout9
\message{[A an endinput that ended the file name]}
\input zzendin\endinput\relax
\message{[B and one that did not]}
\input zzendin \relax
\message{[done]}
\end
